\documentclass[10pt,a4paper]{amsart}
\usepackage[margin=19mm]{geometry}
\usepackage{amsmath,amssymb,mathtools}
\usepackage[scr=rsfs,cal=euler]{mathalfa}
\usepackage{xfrac}
\usepackage[unicode,hidelinks,bookmarks=false]{hyperref}
\newcommand{\ie}{i.e.\ }
\newcommand{\cf}{cf.\ }
\newcommand{\resp}{respectively\ }
\newcommand{\Subsection}{\subsection}
\providecommand{\nobreakdash}{\nobreak\hbox{-}}
\setlength{\emergencystretch}{2em}
\multlinegap0pt
\usepackage{bm}
\usepackage{bbm}
\usepackage{dsfont}
\usepackage{xspace}
\usepackage{cancel}
\usepackage{mathtools}
\usepackage{amsthm}
\makeatletter
\@ifundefined{definition}{\newtheorem{definition}{Definition}}{}
\@ifundefined{lemma}{\newtheorem{lemma}{Lemma}}{}
\@ifundefined{proposition}{\newtheorem{proposition}{Proposition}}{}
\@ifundefined{theorem}{\newtheorem{theorem}{Theorem}}{}
\makeatother
\makeatletter
\expandafter\ifx\csname enbibliography\endcsname\relax
\newenvironment{enbibliography}{\renewcommand{\refname}{{\large\sc References}}\vspace{-0.5cm}\begin{thebibliography}}{\end{thebibliography}}
\fi
\makeatother
\title[Oil displacement by slug injection: a rigorous justification]{Oil displacement by slug injection: a rigorous justification for the Jouguet principle heuristic}
\author{Sergey Matveenko, Nikita Rastegaev}
\date{Source: arXiv:2511.08533v1 (CC BY 4.0)}
\begin{document}
\maketitle

\noindent\emph{Source and licence.} Oil displacement by slug injection: a rigorous justification for the Jouguet principle heuristic. Version arXiv:2511.08533v1 (CC BY 4.0), distributed under the Creative Commons Attribution 4.0 International licence, as recorded in the arXiv licence metadata for this exact version and shown on the abstract page https://arxiv.org/abs/2511.08533v1. Full rights notice: licenses/arxiv-2511.08533-CC-BY-4.0.txt.\par\smallskip

\noindent\textit{Editorial scope.} Counted unit: the Theorem whose source label is \texttt{thm1} in the article (source0.tex lines 1241--1244, source label \texttt{thm1}), delivered together with the article's own complete original proof of it (source0.tex lines 1245--1280, one \texttt{proof} environment of 36 lines). No mathematics is written, repaired or re-derived: statement and proof are the article's own text line for line. Required context retained uncounted: eq:dPhi\_dx\_at\_A (lines 799--802); eq:Joguet\_cond (lines 883--889); eq:U-char-of-c (lines 999--1002); eq:d\_psi\_d\_zeta (lines 1035--1038); eq:psi\_Phi\_definition (lines 1041--1047); def:mapping\_and\_solution\_shorthands (lines 1080--1088); eq:mapping\_into\_zeta\_U (lines 1083--1086); prop:continuous\_differential\_by\_parameter (lines 1090--1093); lemma:dpsi\_dzeta0 (lines 1136--1140); lemma:dpsi\_dU0 (lines 1176--1179); prop:full\_cover\_of\_zeta\_U\_plane (lines 1199--1205). Every definition, display and label of those lines is retained unchanged. Omitted from this package and not redistributed: 1--798, 803--882, 890--998, 1003--1034, 1039--1040, 1048--1079, 1087--1089, 1094--1135, 1141--1175, 1180--1198, 1206--1240, 1281--1527. Nothing inside a retained excerpt is omitted or altered: every retained body is delivered exactly as the article's lines are written, so no omission receipt of any kind accompanies this package. Citations: the retained excerpts carry no citation command at all, so no bibliography entry of the article is transcribed and no reference is redistributed. Reference adaptation: none was needed and none was made; every reference inside the delivered unit resolves inside it, so no unresolved reference or citation is produced. Printed slips kept literal: no printed word, symbol, hypothesis or number of the counted statement or of its proof was altered. Layout and class: the article's own class is \texttt{article} with packages including babel, inputenc, fontenc, hyperref, amssymb, amsmath, amsthm, amsfonts, graphicx, geometry, appendix, cite, float, xcolor; the wrapper is the lane's pinned \texttt{amsart} editorial wrapper at 10pt on A4, which restates the theorem-like environments the retained text needs and adds the article's own macro definitions that the retained text needs (none), with extra wrapper packages (bm, bbm, dsfont, xspace, cancel, mathtools). The delivered numbering is the unit's own. Assets: no retained span contains a figure environment, a graphics-inclusion command, a PDF-page inclusion, a table or a TikZ picture; no image file is copied into this package, the package declares no asset dependency and no image file is redistributed with it. Licence: version arXiv:2511.08533v1 of this article is distributed under the Creative Commons Attribution 4.0 International licence, as recorded in the arXiv licence metadata for this exact version and shown on the abstract page https://arxiv.org/abs/2511.08533v1. A full rights notice is delivered at licenses/arxiv-2511.08533-CC-BY-4.0.txt. Characters: every retained excerpt is pure ASCII, so the delivered engine, XeLaTeX with its default Latin Modern text font, reports no missing character.
\begin{equation}
\label{eq:dPhi_dx_at_A}
\frac{d\Phi}{d x}(x_A) = v(1, 0).
\end{equation}

\begin{definition}
The initial value $\mathcal U$ on $\Phi$ obtained via the \emph{Joguet condition} satisfies the equation
\begin{equation}
\label{eq:Joguet_cond} 
\mathcal F_{\mathcal U}(\mathcal U(\Phi(x), x), \zeta(\Phi(x), x)) = \frac{a(\zeta(\Phi(x), x))}{\zeta(\Phi(x), x)}.
\end{equation} 
\end{definition}

\begin{equation}
\label{eq:U-char-of-c}
\dfrac{d}{d\zeta} \mathcal U = -\dfrac{\mathcal F_{\zeta}(\mathcal U, \zeta)}{\mathcal F_{\mathcal U}(\mathcal U, \zeta) - a_\zeta(\zeta)}.  
\end{equation}

\begin{equation}
\label{eq:d_psi_d_zeta}
    \frac{d\psi}{d\zeta} = \frac{a_{\zeta\zeta}(\zeta)a_\zeta(\zeta)}{1 + a_\zeta^2(\zeta)} + \frac{a_{\zeta\zeta}(\zeta)}{\mathcal F_{\mathcal U}(\mathcal U, \zeta) - a_\zeta(\zeta)}
\end{equation}

\begin{align}
\label{eq:psi_Phi_definition}
\begin{split}
\psi_\Phi(\zeta) &= \frac12\ln((\Phi(x) - t_{inj})^2 + x^2) = \ln x + \dfrac{1}{2}\ln(1+ a^2_\zeta(\zeta_\Phi(x))) \\
{} &= \ln t_{inj} - \ln p(\zeta) + \dfrac{1}{2}\ln(1+ a^2_\zeta(\zeta)).
\end{split}
\end{align}

\begin{definition}
\label{def:mapping_and_solution_shorthands}
For all $\mathcal U$-characteristics in $\angle_{t_{inj}}$ we define the mapping
\begin{equation}
\label{eq:mapping_into_zeta_U}
(x, \varphi(x), \mathcal U(\varphi(x), x)) \to (\zeta, \mathcal U(\zeta), \psi(\zeta; \mathcal U(\cdot))).
\end{equation}
Let $\mathcal U(\zeta;\zeta_0, \mathcal U_0)$ be the solution of~\eqref{eq:U-char-of-c} with initial condition $\mathcal U(\zeta_0) = \mathcal U_0$. For brevity we will use the shorthands $\mathcal U(\zeta;\zeta_0) = \mathcal U(\zeta;\zeta_0, \mathcal U_J(\zeta_0))$ and $\mathcal U(\zeta;\mathcal U_0) = \mathcal U(\zeta;1, \mathcal U_0)$ when appropriate. Similarly, given the solution $\mathcal U(\zeta;\zeta_0)$, we shorthand $\psi(\zeta; \zeta_0) = \psi(\zeta; \mathcal U(\cdot;\zeta_0))$ and for $\mathcal U(\zeta;\mathcal U_0)$, we write $\psi(\zeta;\mathcal U_0) = \psi(\zeta; \mathcal U(\cdot;\mathcal U_0))$.
\end{definition}

\begin{proposition}[Smooth dependence on parameters]
\label{prop:continuous_differential_by_parameter}
The functions $\mathcal U(\zeta;\zeta_0)$, $\mathcal U(\zeta;\mathcal U_0)$, $\psi(\zeta; \zeta_0)$ and $\psi(\zeta;\mathcal U_0)$ are continuously differentiable with respect to their parameters.
\end{proposition}

\begin{lemma} 
\label{lemma:dpsi_dzeta0}
Let ($\mathcal F$5) hold for $\zeta_0$. Then 
$\dfrac{\partial}{\partial\zeta_0}\psi(\zeta;\zeta_0) < 0$ for $\zeta, \zeta_0 \in [0,1]$, $\zeta < \zeta_0$.
\end{lemma}

\begin{lemma}
\label{lemma:dpsi_dU0}
$\dfrac{\partial}{\partial \mathcal U_0}\psi(\zeta; \mathcal U_0) > 0$  for $\zeta\in[0,1]$, $\mathcal U_0\in[1, \mathcal U^+_{OA}]$.
\end{lemma}

\begin{proposition}
\label{prop:full_cover_of_zeta_U_plane}
Let the condition ($\mathcal F$5) hold true for $\zeta(\Phi(x), x)$ at every point $(\Phi(x), x)$  of the chemical shock front. Then the images of $\mathcal U$-characteristics in $\angle_{t_{inj}}$ (continued from TA and constructed via the Jouguet condition \eqref{eq:Joguet_cond} from each point of the chemical shock front $(\Phi(x),x)$) when mapped onto the $(\zeta, \mathcal U)$ plane fill the whole area
$$\overline {\Upsilon}_J = \{(\zeta, \mathcal U): \zeta\in[0,1], \mathcal U \in [1, \mathcal U_J(\zeta)]\}$$
under the graph of $\mathcal U_J$.
Each of $\mathcal U$-characteristic images starts at the segment $\zeta = 0$, $\mathcal U \in [0,\mathcal U_J(0)]$ and arrives either at $\Upsilon_J$ or at $\Upsilon_{TA}$.    
\end{proposition}

\begin{theorem}
\label{thm1}
Let the condition ($\mathcal F$5) hold true for $\zeta(\Phi(x), x)$ at every point $(\Phi(x), x)$  of the chemical shock front. Then the $\mathcal U$-characteristics in $\angle_{t_{inj}}$ (continued from TA and constructed via the Jouguet condition \eqref{eq:Joguet_cond} from each point of the chemical shock front $(\Phi(x),x)$) fill the whole area $\angle_{t_{inj}}$ without intersections and give us a pice-wise continuously differentiable solution $\mathcal U(\varphi, x)$.
\end{theorem}

\begin{proof}
Note that the characteristic from the point $A$ as the continuation of the characteristic $OA$, and the characteristic from $A$ constructed via the Jouguet condition give us the same curve due to \eqref{eq:dPhi_dx_at_A}. Therefore we have two families of characteristics that have the $A$ characteristic in common.

Similar to Proposition~\ref{prop:full_cover_of_zeta_U_plane} we first map all characteristics into the $(\zeta, \mathcal U)$ plane using the mapping \eqref{eq:mapping_into_zeta_U}. In the $(\zeta, \mathcal U)$ plane the curves $\mathcal U(\zeta; \mathcal U_0)$ correspond to the characteristics continued from TA, while the curves $\mathcal U(\zeta; \zeta_0)$ correspond to the Jouguet characteristics (see Definition~\ref{def:mapping_and_solution_shorthands} for the solution shorthands we use here).

Due to Proposition~\ref{prop:full_cover_of_zeta_U_plane} together the curves $\mathcal U(\zeta; \mathcal U_0)$ and $\mathcal U(\zeta; \zeta_0)$ fill the whole area $\overline{\Upsilon}_J$.

Now, we consider the solutions $\psi(\zeta; \mathcal U_0)$ and $\psi(\zeta; \zeta_0)$ and their inverse mapping
\[
x = \frac{e^\psi}{\sqrt{a_\zeta^2(\zeta) + 1}}, \qquad \varphi = t_{inj} + \frac{e^{\psi}a_\zeta(\zeta)}{\sqrt{a_\zeta^2(\zeta) + 1}}
\]
back into the $(\varphi, x)$ plane. Due to Proposition~\ref{prop:continuous_differential_by_parameter}, Lemma~\ref{lemma:dpsi_dzeta0} and Lemma~\ref{lemma:dpsi_dU0} they are continuous and monotone with respect to their parameters. On the upper boundary $\Upsilon_J$ the functions $\psi(\zeta; \zeta_0)$ have boundary values $\psi(\zeta_0; \zeta_0) = \psi_\Phi(\zeta_0)$, so when we go back to $(\varphi, x)$ coordinates, this boundary will map back into the shock front $\Phi$. Similarly, on the right boundary of $\overline{\Upsilon}_J$ we have $\psi(1; \mathcal U_0) = \psi_{TA}(\mathcal U_0)$, which will map back into $TA$. On the lower boundary we have $\mathcal U \equiv 1$, so when approaching that boundary, $\psi(1; \mathcal U_0) \to -\infty$ as $\mathcal U_0\to 1$, and the right-hand side of \eqref{eq:d_psi_d_zeta} is bounded, so $\psi(\zeta; \mathcal U_0) \to -\infty$ as $\mathcal U_0\to 1$. Therefore, the lower boundary will map back into the point $T$. Finally, the left boundary maps into the line $\varphi = t_{inj}+a_\zeta(0)x$, since $\zeta = 0$ on it, and as $\mathcal U$ approaches $\mathcal U_J(0)$, the values of $\psi(\zeta; \zeta_0)$ must go to infinity, since $\psi_\Phi(\zeta) \to \infty$ as $\zeta\to 0$ by the definition \eqref{eq:psi_Phi_definition}.
In summary, we demonstrated that the boundaries of $\overline{\Upsilon}_J$ map back into the boundaries of $\angle_{t_{inj}}$, and since $\psi$ is continuous and monotone with respect to the parameters, the interior will also map one-to-one. Thus, the characteristics in the $(\varphi, x)$ plane cover the whole area $\angle_{t_{inj}}$.

What's left is to prove that $\mathcal U$, when mapped back into the $(\varphi, x)$ plane, gives us a piecewise continuously differentiable function. We can prove this separately for the families $\mathcal U(\zeta; \mathcal U_0)$ and $\mathcal U(\zeta; \zeta_0)$. The following proof is for the $\mathcal U(\zeta; \zeta_0)$ family, the proof for the other family is the same.

Denote by
\[
L(\zeta; \zeta_0) = \left(t_{inj} + \frac{e^{\psi(\zeta; \zeta_0)}a_\zeta(\zeta)}{\sqrt{a_\zeta^2(\zeta) + 1}}, \frac{e^{\psi(\zeta; \zeta_0)}}{\sqrt{a_\zeta^2(\zeta) + 1}}\right).
\]
Due to Proposition~\ref{prop:continuous_differential_by_parameter} this is a $\mathcal C^1$ mapping. Using this mapping we can define the inverse mapping of $\mathcal U(\zeta; \zeta_0)$ back into the $(\varphi, x)$ plane as follows:
\[
\mathcal U(\varphi, x) := \mathcal U(L^{-1}(\varphi, x)).
\]
For the inverse mapping $L^{-1}$ to exist and be continuously differential, we need its Jacobian determinant to be non-zero. Denote
\[
G(\zeta, \zeta_0) = \frac{e^{\psi(\zeta; \zeta_0)}}{\sqrt{a_\zeta^2(\zeta) + 1}}, \qquad
L(\zeta, \zeta_0) = \Big(t_{inj} + a_\zeta(\zeta)G(\zeta, \zeta_0), G(\zeta, \zeta_0)\Big).
\]
Then the Jacobian determinant is calculated as follows:
\begin{align*}
\det J_L &= (a_{\zeta\zeta}(\zeta) G(\zeta, \zeta_0) + a_{\zeta} G_\zeta(\zeta, \zeta_0)) G_{\zeta_0}(\zeta, \zeta_0) - a_\zeta(\zeta) G_{\zeta_0}(\zeta, \zeta_0) G_\zeta(\zeta, \zeta_0) \\
{} &= a_{\zeta\zeta}(\zeta) G(\zeta, \zeta_0) G_{\zeta_0}(\zeta, \zeta_0) = a_{\zeta\zeta}(\zeta) G^2(\zeta, \zeta_0) \psi_{\zeta_0}(\zeta; \zeta_0) > 0
\end{align*}
for $\zeta<\zeta_0$ due to Lemma~\ref{lemma:dpsi_dzeta0}.
\end{proof}

\end{document}
